\documentclass[11pt]{article}
\usepackage[margin=1in]{geometry}
\usepackage[T1]{fontenc}
\usepackage[utf8]{inputenc}
\usepackage{lmodern,microtype}
\usepackage{amsmath,amssymb,booktabs,longtable,array}
\usepackage{xcolor,hyperref}
\hypersetup{
  colorlinks=true,
  linkcolor=blue!45!black,
  urlcolor=blue!45!black,
  pdftitle={Resolution Theory: Complete Foundations, Proof Paths, and Beyond the Finite},
  pdfauthor={Komorebi P.},
  pdfcreator={},
  pdfproducer={}
}
\setlength{\parindent}{0pt}
\setlength{\parskip}{7pt}
\newcolumntype{L}[1]{>{\raggedright\arraybackslash}p{#1}}
\begin{document}
\thispagestyle{empty}
\begin{center}
\vspace*{1.0in}
{\fontsize{30}{34}\selectfont\bfseries Resolution Theory\par}
\vspace{0.22in}
{\Large Complete Foundations, Proof Paths, and Beyond the Finite\par}
\vspace{0.30in}
{\large A collected work on representation, refinement, nonlinear interaction, logic, compression, categorical realization, and infinite resolution\par}
\vfill
{\Large Komorebi P.\par}
\vspace{0.75in}
\end{center}
\clearpage

\section*{Nothing has been replaced}
This collected edition restores the structure that the shorter companion obscured.  The finite theory, nonlinear theory, Resolution Logic, universal representation theorem, proof-path methodology, and final philosophical synthesis remain intact.  \emph{Beyond the Finite} follows them as a second movement: it deepens the countable, unbounded, measurable, and categorical side of the theory without standing in for the rest.

The basic object of the theory is not an isolated mathematical object.  It is an object together with the resolutions through which it can be seen, the transfers and nonlinear interactions that act between its packets, the certificates that record approximation and cost, and the proof paths by which its claims became visible.

\[
\mathfrak P
=
(X,\mathcal R,\mathcal T,\mathcal J,\mathcal C,\mathcal E,\mathcal A).
\]

Here $X$ is the object or problem, $\mathcal R$ its admissible resolutions, $\mathcal T$ its linear transfer field, $\mathcal J$ its nonlinear interaction hypergraph, $\mathcal C$ its quantitative certificate field, $\mathcal E$ its typed evidence, and $\mathcal A$ its expanded archive of proofs, examples, failed routes, and unresolved questions.

\section*{How the collected work is arranged}
\begin{longtable}{L{1.05in}L{2.25in}L{3.65in}}
\toprule
Division & Main question & Contents\\
\midrule
\endfirsthead
\toprule
Division & Main question & Contents\\
\midrule
\endhead
Part I & What is a resolution-aware mathematical object? & Finite projector systems, transfer corners, refinement, compression, entropy, atlases, gluing, and resolution complexity.\\
Part II & Can every abstract resolved system be represented concretely? & Linking algebras, faithful projector-corner representations, universal properties, recognition, and minimal realizations.\\
Part III & What survives beyond finite dimension? & Countable resolutions, strong and norm topologies, unbounded operators, forms, resolvents, normal joins, measurable packets, and categorical completion.\\
Part IV & What happens when the maps are nonlinear? & Resolved jets, hyperedges, interaction trees, Fa\`a di Bruno composition, nonlinear inversion, dynamics, memory, and stochastic corrections.\\
Part V & What can be asserted at a stated resolution? & Resolution Logic, quantitative judgments, evidence types, refinement and coarsening rules, gluing, and soundness.\\
Part VI & What is lost when mathematics is compressed into a final proof? & Proof paths, compressed certificates, bridge value, collective noise, provenance, audit capacity, and representation portfolios.\\
Part VII & What do mathematics, computer science, statistics, and technology mean inside one architecture? & The final synthesis and the enlarged closure theorem for resolution-aware practice.\\
Part VIII & Where does the finite picture stop being honest? & The complete companion \emph{Beyond the Finite}: compactness boundaries, trace-class states, graph cores, strong-resolvent limits, multiplier topology, direct integrals, measurable descent, and practical proof workflows.\\
\bottomrule
\end{longtable}

\section*{Status of the mathematics}
The collected work separates four kinds of statement throughout:
\[
\boxed{
\text{proved}
\quad\text{conditional}
\quad\text{computationally verified}
\quad\text{open}
}
\]
A proof is not upgraded by a numerical experiment, and a conjecture is not hidden inside an implementation.  Conversely, an exact computation is retained as exact evidence rather than being dismissed because it is not a symbolic proof.  The point is not to flatten every kind of knowledge into one label, but to keep the interfaces between them honest.

\section*{Reading principle}
The first volume should be read as the complete theory.  The companion should be read as a sustained examination of one boundary of that theory.  Its purpose is not to repeat the finite results, but to show what changes when sums become countable, units move to multiplier algebras, operators acquire domains, and finite calculations need topology-specific certificates.

No section in the companion cancels or replaces the proof-path, nonlinear, logical, or philosophical parts of the main theory.  The two works belong together precisely because they answer different questions.
\end{document}
